\documentclass[aspectratio=169,10pt]{beamer}
\usetheme[numbering=fraction,progressbar=frametitle]{metropolis}
\usepackage{fontspec}
\setsansfont{Latin Modern Sans}
\usepackage{amsmath,amssymb,mathtools,graphicx,booktabs}
\usepackage{polyglossia}
\setdefaultlanguage{english}
\setbeamertemplate{navigation symbols}{}
\setbeamerfont{frametitle}{size=\large}
\setbeamerfont{normal text}{size=\normalsize}
\setlength{\parskip}{.45em}
\hypersetup{pdfauthor={},pdfsubject={Lecture slides},pdfkeywords={},pdftitle={Consolidation and transition}}
\providecommand{\N}{\mathbb{N}}
\providecommand{\Z}{\mathbb{Z}}
\providecommand{\D}{\mathbb{D}}
\providecommand{\Q}{\mathbb{Q}}
\providecommand{\R}{\mathbb{R}}
\providecommand{\C}{\mathbb{C}}
\providecommand{\vect}[1]{\overrightarrow{#1}}
\providecommand{\vi}{\vec{\imath}}
\providecommand{\vj}{\vec{\jmath}}
\providecommand{\vk}{\vec{k}}
\providecommand{\norm}[1]{\left\lVert #1\right\rVert}
\providecommand{\vc}[2]{\begin{pmatrix}#1\\#2\end{pmatrix}}
\providecommand{\deter}[2]{\det\left(#1,#2\right)}
\providecommand{\abs}[1]{\left\lvert #1\right\rvert}
\providecommand{\dd}{\,\mathrm{d}}
\providecommand{\e}{\mathrm{e}}
\providecommand{\Card}{\operatorname{Card}}
\providecommand{\E}{\operatorname{E}}
\providecommand{\V}{\operatorname{V}}
\providecommand{\degre}{^{\circ}}
\providecommand{\un}[1]{\,\mathrm{#1}}
\providecommand{\sci}[2]{#1\times10^{#2}}
\providecommand{\pgcd}{\operatorname{PGCD}}
\providecommand{\ppcm}{\operatorname{PPCM}}
\title{Consolidation and transition}
\subtitle{Lecture and worked examples}
\date{}
\author{}
\begin{document}
\begin{frame}[plain]
\vfill
{\usebeamerfont{title}\inserttitle\par}
\medskip
{\usebeamerfont{subtitle}\insertsubtitle\par}
\medskip
{\usebeamercolor[fg]{progress bar}\rule{\textwidth}{.5pt}}
\vfill
\end{frame}
\begin{frame}{Review of calculation, functions and geometry}
\begin{enumerate}
\setcounter{enumi}{0}
\item Reduce \(168/252\) to lowest terms; calculate \(2^{-3}+1/4\).

\item Factorise \(x^2-6x+9\), then solve \((x-3)^2=4\).

\item Solve \(\dfrac{x-2}{x+1}\geq0\).

\item For \(f(x)=4-(x-1)^2\) on \([-2,3]\), give the maximum and the abscissae where \(f(x)=0\).

\item With \(A(0,1)\), \(B(4,3)\), \(C(6,4)\), prove collinearity and give an equation of the line.

\end{enumerate}
\end{frame}
\begin{frame}{Worked answer: Review of calculation, functions and geometry}
1. \(168/252=2/3\), and \(2^{-3}+1/4=1/8+2/8=3/8\). 2. \((x-3)^2\); \(x-3=\pm2\), so \(S=\{1,5\}\). 3. Forbidden \(-1\), zero \(2\); signs \(+,-,+\), so \(S=(-\infty,-1)\cup[2,+\infty)\). 4. The square is positive or zero, so \(f(x)\leq4\), with equality at \(x=1\). The zeros are \(x=-1,3\), both in the domain. 5. \(\vect{AB}=(4,2)\), \(\vect{AC}=(6,3)=\dfrac32\vect{AB}\). An equation is \(x-2y+2=0\).
\end{frame}
\begin{frame}{Review of data, probability and modelling}
\begin{enumerate}
\setcounter{enumi}{0}
\item Calculate the mean and median of \(2;4;4;6;9\). Determine \(Q_1,Q_3\).

\item A price drops by \(15\%\), then rises by \(20\%\). Give the overall rate and the rate that then restores the initial price.

\item In a fictional population of \(100\) people, \(40\) play sport (\(S\)), \(30\) come by bike (\(V\)), \(18\) do both. Construct the two-way table and calculate \(P_S(V)\), \(P_V(S)\) for an equally likely draw.

\item For a rectangle of perimeter \(24\) m, express the area as a function of one dimension \(x\) and prove its maximum.

\end{enumerate}
\end{frame}
\begin{frame}{Worked answer: Review of data, probability and modelling}
1. Mean \(25/5=5\), median \(4\). Quartile ranks \(2\) and \(4\), so \(Q_1=4\), \(Q_3=6\).

2. Overall multiplier \(0.85\times1.20=1.02\): an increase of \(2\%\). Reverse rate \(1/1.02-1=-1/51\approx-1.96\%\).

3. Table: \(S\cap V=18\), \(S\cap\overline V=22\), \(\overline S\cap V=12\), \(\overline S\cap\overline V=48\). \(P_S(V)=18/40=45\%\), \(P_V(S)=18/30=60\%\).

4. \(A(x)=x(12-x)=36-(x-6)^2\) on \((0,12)\). The maximum is \(36\) m\(^2\), attained for the square of side \(6\) m.
\end{frame}
\begin{frame}{Choosing a strategy and preparing what comes next}
For a given function \(f\), explain the following program. What does it return with \(f(x)=x^2\)?
\end{frame}
\begin{frame}{Worked answer: Choosing a strategy and preparing what comes next}
It returns a list of abscissa-image pairs for the integers from \(-3\) to \(3\) inclusive: \((-3,9),(-2,4),(-1,1),(0,0),(1,1),(2,4),(3,9)\). The bound \(4\) of \texttt{range} is excluded. The table does not describe the images of all real numbers between \(-3\) and \(3\).
\end{frame}
\begin{frame}[fragile]{Python}
\small
\begin{verbatim}
def tableau(f):
    valeurs = []
    for x in range(-3, 4):
        valeurs.append((x, f(x)))
    return valeurs
\end{verbatim}
\end{frame}
\begin{frame}{Choosing a strategy and preparing what comes next}
\textbf{Search}: try things out, organise leads. \textbf{Model}: choose variables, domain and assumptions. \textbf{Represent}: move from text to table, figure or graph. \textbf{Reason}: justify and distinguish conjecture from proof. \textbf{Calculate}: choose a form and check the operations. \textbf{Communicate}: explain clearly with suitable notation and units.
\end{frame}
\begin{frame}{Choosing a strategy and preparing what comes next}
Choose a successful piece of work and a corrected one. For each, identify the skills used, explain one method and state a personal goal for the following year. Then redo an exercise matching this goal without looking at the correction.
\end{frame}
\begin{frame}{Worked answer: Choosing a strategy and preparing what comes next}
There is no single numerical answer. A satisfactory piece of work precisely identifies the problem, the chosen method, a justification and a check. The goal must be observable, for example: “state the forbidden values before solving a quotient equation”, rather than “get better at calculation”.
\end{frame}
\end{document}
